\documentclass[11pt,a4paper]{article}
\usepackage[utf8]{inputenc}\usepackage[T1]{fontenc}\usepackage{mathptmx}
\usepackage[margin=2.2cm]{geometry}\usepackage{booktabs}\usepackage{tabularx}
\usepackage{array}\usepackage{xcolor}\usepackage{titlesec}\usepackage{enumitem}
\usepackage{hyperref}\usepackage{fancyhdr}\usepackage{setspace}\usepackage{parskip}
\usepackage{float}\usepackage{colortbl}

\definecolor{navy}{HTML}{1B2A4A}\definecolor{gold}{HTML}{C49A2A}
\definecolor{dk}{HTML}{2D2D2D}\definecolor{mg}{HTML}{4A4A4A}
\definecolor{sb}{HTML}{D6E4F0}\definecolor{sg}{HTML}{DFF0D8}
\definecolor{so}{HTML}{FDEBD0}\definecolor{sp}{HTML}{E8DAEF}
\definecolor{sr}{HTML}{FADBD8}

\definecolor{revcol}{HTML}{2E5090}\definecolor{rdcol}{HTML}{2E7D32}
\definecolor{forcol}{HTML}{E65100}\definecolor{hlcol}{HTML}{6A1B9A}
\definecolor{revbg}{HTML}{D6E4F0}\definecolor{rdbg}{HTML}{DFF0D8}
\definecolor{forbg}{HTML}{FDEBD0}\definecolor{hlbg}{HTML}{E8DAEF}

\titleformat{\section}{\Large\bfseries\color{navy}}{\thesection}{1em}{}[\vspace{-4pt}{\color{gold}\rule{\textwidth}{1.5pt}}]
\titleformat{\subsection}{\large\bfseries\color{navy}}{\thesubsection}{1em}{}
\setlength{\headheight}{14pt}\pagestyle{fancy}\fancyhf{}
\fancyhead[L]{\small\color{mg}\textit{AIGGPA --- Department-Specific Schedule}}
\fancyhead[R]{\small\color{mg}\textit{April 2026}}
\fancyfoot[C]{\small\color{mg}\thepage}
\renewcommand{\headrulewidth}{0.4pt}
\hypersetup{colorlinks=true,linkcolor=navy,citecolor=navy,urlcolor=navy}
\setstretch{1.1}
\newcolumntype{L}[1]{>{\raggedright\arraybackslash}p{#1}}
\newcolumntype{C}[1]{>{\centering\arraybackslash}p{#1}}

\newcommand{\qcard}[8]{%
\noindent\begin{tabularx}{\textwidth}{@{}>{\columncolor{#7}\bfseries\small}L{3.2cm}|X@{}}
\arrayrulecolor{#8}\hline
Theoretical Construct & #1 \\
Dimension & #2 \\
Operational Indicator & #3 \\
\arrayrulecolor{gold}\hline
\rowcolor{white}\textbf{\color{#8}Survey Item} & \textbf{\color{#8}#4} \\
\arrayrulecolor{gold}\hline
Data Type \& Scale & #5 \\
Justification & \textit{#6} \\
\arrayrulecolor{#8}\hline
\end{tabularx}\vspace{10pt}
}

\begin{document}
\begin{titlepage}\centering
\vspace*{2.5cm}
{\color{gold}\rule{0.6\textwidth}{2pt}}\vspace{1cm}
{\Huge\bfseries\color{navy} DEPARTMENT-SPECIFIC\\[6pt] QUESTIONNAIRE SUPPLEMENT}\vspace{0.5cm}
{\Large\color{dk} Structured Schedule Items Tailored to\\[4pt] Revenue, Rural Development, Forest \& Health}\vspace{1cm}
{\color{gold}\rule{0.6\textwidth}{2pt}}\vspace{2cm}
{\large\color{mg} This document supplements the main questionnaire (39 items)\\[4pt] with 20 department-specific questions --- 5 per department ---\\[4pt] targeting the unique digital ecosystems of each department.}\vspace{2cm}
\begin{tabular}{r l}
\textbf{\color{navy}Prepared By:} & Aryan Kori \\[6pt]
\textbf{\color{navy}Institution:} & AIGGPA, Bhopal \\[6pt]
\textbf{\color{navy}Date:} & April 2026 \\
\end{tabular}
\vfill
{\small\color{mg} Administer only the section relevant to the respondent's department.}
\end{titlepage}
\tableofcontents\newpage

\section*{How to Use This Supplement}
\addcontentsline{toc}{section}{How to Use This Supplement}
This document contains \textbf{20 department-specific questions} (5 per department) designed to complement the 39-item main questionnaire. Each respondent receives \textbf{only the section matching their department}, resulting in a total instrument of 39~+~5~=~\textbf{44 questions per respondent}.

\begin{table}[H]\centering\small
\begin{tabularx}{\textwidth}{L{3.5cm} L{3cm} X}
\toprule
\textbf{Department} & \textbf{Questions} & \textbf{Key Digital Ecosystem} \\
\midrule
\textcolor{revcol}{\textbf{Revenue}} & Q40--Q44 & Bhulekh/WebGIS, RCMS, SAARA, SAMPADA \\[4pt]
\textcolor{rdcol}{\textbf{Rural Development}} & Q45--Q49 & NREGASoft/NMMS, e-Gram Swaraj, PMAY-G, SBM-G \\[4pt]
\textcolor{forcol}{\textbf{Forest}} & Q50--Q54 & e-Green Watch, AI Alert System, GIS/Remote Sensing \\[4pt]
\textcolor{hlcol}{\textbf{Health}} & Q55--Q59 & ANMOL MP, HMIS, Nikshay, eVIN, IHIP, ABHA \\
\bottomrule
\end{tabularx}
\end{table}
\newpage

\section{\textcolor{revcol}{Revenue Department}}

\textbf{Core Functions:}

\begin{table}[H]\centering\small
\begin{tabularx}{\textwidth}{L{3cm} X}
\toprule
\textbf{Core Function} & \textbf{Description} \\
\midrule
Land Record Management & Maintenance, mutation, and verification of land ownership records (Khasra, B-1, land maps) \\[3pt]
Revenue Case Adjudication & Hearing and disposal of revenue disputes (demarcation, partition, succession) through Tehsil and SDM courts \\[3pt]
Registration \& Stamps & Registration of property sale/purchase deeds, lease agreements, and collection of stamp duty \\[3pt]
Revenue Collection & Collection of land revenue, cess, and other government dues at tehsil/district level \\[3pt]
Disaster \& Relief & Assessment and disbursement of crop damage compensation and disaster relief under SDRF/NDRF \\
\bottomrule
\end{tabularx}
\end{table}


\textbf{Cadre Structure \& Digital Responsibilities:}

\begin{table}[H]\centering\small
\begin{tabularx}{\textwidth}{L{1.2cm} L{3.5cm} L{4.5cm} X}
\toprule
\textbf{Class} & \textbf{Designations} & \textbf{Core Work} & \textbf{Digital Role} \\
\midrule
\textbf{I} & Collector, Additional Collector, Commissioner (Land Records) & District-level policy, appellate authority, disaster relief coordination, supervision of all revenue operations & Dashboard review, MIS monitoring, approval workflows in SAARA/RCMS; rarely direct data entry \\[4pt]
\textbf{II} & SDM (Sub-Divisional Magistrate), Deputy Collector, Dy. Director (Land Records) & Sub-divisional case adjudication, land acquisition, election duty, magisterial functions & RCMS case disposal, Bhulekh verification for land acquisition, e-Court hearings; moderate data entry \\[4pt]
\textbf{III} & Tehsildar, Naib Tehsildar, Revenue Inspector (RI), Patwari, Computer Operator & Land record maintenance, mutation, demarcation, revenue collection, Khasra/B-1 updates, data entry & Primary users of Bhulekh/WebGIS (Patwari), RCMS (Tehsildar), SAMPADA (registration staff); heaviest digital workload \\[4pt]
\textbf{IV} & Peon, Office Assistant, Process Server, Mali, Chowkidar & File movement, court summons delivery, office upkeep, record room management & Minimal direct use; may assist with scanning, photocopying, or carrying files to digital entry points \\
\bottomrule
\end{tabularx}
\end{table}

\textbf{Key Digital Tools:} MP Bhulekh / WebGIS 2.0, RCMS (Revenue Case Management System), SAARA (Smart Application for Revenue Administration), SAMPADA 2.0 (Registration \& Stamps), e-Court / e-Filing

\subsection*{Q40. Which of the following Revenue department digital tools are you aware of? (Select all...}
\qcard{Facilitating Conditions (FC)}{Department-Specific Tool Awareness}{Awareness of Revenue-specific digital platforms}{Which of the following Revenue department digital tools are you aware of? (Select all that apply)}{Multi-select Checklist (MP Bhulekh/WebGIS 2.0 / RCMS / SAARA / SAMPADA 2.0 / e-Court/e-Filing / None)}{Revenue department operates a distinct digital ecosystem for land records, case management, and registration. Awareness must be measured at the department-specific tool level, not just generic tools, to identify precise FC gaps (Ndou, 2004).}{revbg}{revcol}

\subsection*{Q41. To what extent has the MP Bhulekh/WebGIS 2.0 portal improved the speed and accuracy o...}
\qcard{Performance Expectancy (PE)}{Task-Technology Fit (Land Records)}{Perceived usefulness of Bhulekh/WebGIS for land record management}{To what extent has the MP Bhulekh/WebGIS 2.0 portal improved the speed and accuracy of land record verification in your office?}{5-point Improvement Scale (1=No improvement to 5=Greatly improved)}{Land record digitisation is the flagship digital intervention in Revenue; measuring its perceived impact directly tests PE for the department's core workflow (Davis, 1989).}{revbg}{revcol}

\subsection*{Q42. How difficult do you find the Revenue Case Management System (RCMS) to use for tracki...}
\qcard{Effort Expectancy (EE)}{System Complexity (RCMS)}{Perceived difficulty of using RCMS for case tracking}{How difficult do you find the Revenue Case Management System (RCMS) to use for tracking and managing revenue cases?}{5-point Likert (1=Very Easy to 5=Very Difficult)}{RCMS is a complex workflow system requiring multi-step data entry; high EE here indicates need for role-specific RCMS training rather than generic IT training (Venkatesh et al., 2003).}{revbg}{revcol}

\subsection*{Q43. What percentage of revenue cases or land records in your office still require referen...}
\qcard{Facilitating Conditions (FC)}{Data Migration Challenges}{Extent of legacy paper records still requiring manual processing}{What percentage of revenue cases or land records in your office still require reference to physical (paper) files that have not been digitised?}{Categorical (0-20 percent / 21-40 percent / 41-60 percent / 61-80 percent / 81-100 percent)}{Incomplete digitisation of legacy records is a Revenue-specific FC barrier; even with functional digital tools, employees are forced into hybrid workflows if historical data remains on paper (Heeks, 2006).}{revbg}{revcol}

\subsection*{Q44. How often do citizens or applicants visit your office expecting services that require...}
\qcard{Social Influence (SI)}{Citizen-Facing Pressure}{Frequency of citizen enquiries that require digital tool usage}{How often do citizens or applicants visit your office expecting services that require you to use digital tools (e.g., online land record verification, e-stamp)?}{5-point Frequency (1=Never to 5=Multiple times daily)}{Revenue is a high-volume citizen-facing department; external demand for digital services creates SI pressure that may accelerate adoption independently of internal mandates (Cordella \& Bonina, 2012).}{revbg}{revcol}


\section{\textcolor{rdcol}{Rural Development Department}}

\textbf{Core Functions:}

\begin{table}[H]\centering\small
\begin{tabularx}{\textwidth}{L{3cm} X}
\toprule
\textbf{Core Function} & \textbf{Description} \\
\midrule
Employment Guarantee & Implementation of MGNREGA --- job card issuance, work allocation, muster rolls, and wage disbursement \\[3pt]
Rural Housing & Beneficiary identification, sanctioning, fund transfer, and geo-tagged monitoring under PMAY-Gramin \\[3pt]
Panchayat Governance & Digital planning, accounting, and reporting for Gram Panchayats via e-Gram Swaraj and GPDP preparation \\[3pt]
Sanitation (SBM-G) & Tracking ODF Plus status, toilet construction, solid/liquid waste management at village level \\[3pt]
Rural Infrastructure & Planning and monitoring of roads, bridges, and public buildings through Rural Engineering Services \\
\bottomrule
\end{tabularx}
\end{table}


\textbf{Cadre Structure \& Digital Responsibilities:}

\begin{table}[H]\centering\small
\begin{tabularx}{\textwidth}{L{1.2cm} L{3.5cm} L{4.5cm} X}
\toprule
\textbf{Class} & \textbf{Designations} & \textbf{Core Work} & \textbf{Digital Role} \\
\midrule
\textbf{I} & CEO (Zila Panchayat), Project Director (DRDA), Joint Director & District-level planning, scheme monitoring, fund allocation, convergence across schemes & Dashboard review on e-Gram Swaraj, PFMS fund release approvals, MIS reporting to state; supervisory digital role \\[4pt]
\textbf{II} & BDO/CEO (Janpad Panchayat), Asst. Project Officer, Dy. CEO & Block-level scheme implementation, staff supervision, expenditure monitoring, inspection & NREGASoft work approval, PMAY-G beneficiary sanctioning, e-Gram Swaraj GPDP uploads; regular portal use \\[4pt]
\textbf{III} & Gram Panchayat Secretary, Gram Rozgar Sahayak (GRS), Accountant, Data Entry Operator, Gram Sevak & Muster roll maintenance, wage list preparation, beneficiary registration, village-level record keeping & Heaviest users --- NREGASoft data entry, NMMS attendance uploads, PMAY-G geo-tagging, SBM-G reporting; daily multi-portal operations \\[4pt]
\textbf{IV} & Peon, Office Assistant, Chowkidar & File movement, office maintenance, physical delivery of notices & Minimal direct use; may carry printed MIS reports or assist with hardware maintenance \\
\bottomrule
\end{tabularx}
\end{table}

\textbf{Key Digital Tools:} NREGASoft MIS / NMMS App (MGNREGA), e-Gram Swaraj, PMAY-G Portal, SBM-G Portal, Panchayat Darpan, PFMS

\subsection*{Q45. Which of the following Rural Development department digital tools are you aware of? (...}
\qcard{Facilitating Conditions (FC)}{Department-Specific Tool Awareness}{Awareness of Rural Development digital platforms}{Which of the following Rural Development department digital tools are you aware of? (Select all that apply)}{Multi-select Checklist (NREGASoft/NMMS App / e-Gram Swaraj / PMAY-G Portal / SBM-G Portal / Panchayat Darpan / PFMS / None)}{Rural Development operates the highest number of scheme-specific MIS portals; awareness fragmentation across multiple platforms is a unique FC challenge for this department.}{rdbg}{rdcol}

\subsection*{Q46. How difficult do you find it to manage and update data across multiple portals (e.g.,...}
\qcard{Effort Expectancy (EE)}{Multi-Portal Burden}{Perceived difficulty of managing multiple scheme portals simultaneously}{How difficult do you find it to manage and update data across multiple portals (e.g., NREGASoft, e-Gram Swaraj, PMAY-G) simultaneously?}{5-point Likert (1=Very Easy to 5=Very Difficult)}{Unlike other departments, Rural Development staff must operate 4-6 separate MIS platforms daily; this multi-portal burden is a department-specific EE barrier not captured by generic difficulty questions (Venkatesh et al., 2003).}{rdbg}{rdcol}

\subsection*{Q47. When working at block-level or panchayat-level offices, how would you rate the intern...}
\qcard{Facilitating Conditions (FC)}{Field-Level Connectivity}{Reliability of internet access at block/panchayat level offices}{When working at block-level or panchayat-level offices, how would you rate the internet connectivity available for uploading data to portals?}{5-point Likert (1=No connectivity to 5=Excellent)}{Rural Development uniquely operates at the panchayat/block level where connectivity is significantly worse than district HQs; this tests FC at the actual operational level rather than the office level (World Bank, 2016).}{rdbg}{rdcol}

\subsection*{Q48. To what extent has the NMMS (National Mobile Monitoring System) app improved the accu...}
\qcard{Performance Expectancy (PE)}{Geo-Tagging \& Mobile Monitoring}{Perceived usefulness of NMMS app for field verification}{To what extent has the NMMS (National Mobile Monitoring System) app improved the accuracy of attendance and worksite verification under MGNREGA?}{5-point Improvement Scale (1=No improvement to 5=Greatly improved)}{NMMS mandates real-time geo-tagged attendance; testing its perceived impact measures PE for mobile-first field monitoring tools specific to rural schemes (Government of India, 2020).}{rdbg}{rdcol}

\subsection*{Q49. On average, what proportion of your working day is spent on entering data into variou...}
\qcard{Effort Expectancy (EE)}{Data Entry Burden}{Time spent on digital data entry vs. fieldwork}{On average, what proportion of your working day is spent on entering data into various portals and MIS systems?}{Categorical (Less than 1 hour / 1-2 hours / 2-4 hours / More than 4 hours / Almost the entire day)}{A known complaint in Rural Development is that portal compliance consumes time meant for field implementation; quantifying this tests whether digital tools enhance or impede actual service delivery (Heeks, 2003).}{rdbg}{rdcol}


\section{\textcolor{forcol}{Forest Department}}

\textbf{Core Functions:}

\begin{table}[H]\centering\small
\begin{tabularx}{\textwidth}{L{3cm} X}
\toprule
\textbf{Core Function} & \textbf{Description} \\
\midrule
Forest Protection & Patrolling, detection and prosecution of illegal felling, encroachment, poaching, and forest fires \\[3pt]
Wildlife Conservation & Monitoring of tiger reserves, national parks, and wildlife corridors; anti-poaching operations \\[3pt]
Afforestation (CAMPA) & Planning, execution, and geo-tagged monitoring of compensatory afforestation and plantation drives \\[3pt]
Forest Inventory & Maintaining updated records of forest area, tree density, species distribution using GIS/remote sensing \\[3pt]
Community Forestry & Joint Forest Management (JFM), tribal rights under FRA, and eco-tourism management \\
\bottomrule
\end{tabularx}
\end{table}


\textbf{Cadre Structure \& Digital Responsibilities:}

\begin{table}[H]\centering\small
\begin{tabularx}{\textwidth}{L{1.2cm} L{3.5cm} L{4.5cm} X}
\toprule
\textbf{Class} & \textbf{Designations} & \textbf{Core Work} & \textbf{Digital Role} \\
\midrule
\textbf{I} & CCF (Chief Conservator), CF (Conservator), DFO/DCF (Divisional Forest Officer) --- IFS cadre & Division/circle management, CAMPA fund utilisation, wildlife protection policy, working plan approval & e-Green Watch dashboard, AI alert system monitoring, GIS-based planning; approval and review role \\[4pt]
\textbf{II} & ACF (Assistant Conservator), Sub-Divisional Forest Officer (SDFO) --- SFS cadre & Sub-division supervision, offence prosecution, nursery inspection, plantation monitoring & Forest Offence MIS entries, e-Green Watch geo-tagging verification, GIS map review; moderate digital use \\[4pt]
\textbf{III} & Range Officer (RFO/Ranger), Deputy Ranger, Forester & Range-level patrolling, offence detection, fire management, plantation execution, wildlife census fieldwork & AI alert response (field verification with GPS photos), GIS data collection, Forest Offence MIS first entry; field-heavy digital use \\[4pt]
\textbf{IV} & Forest Guard, Van Rakshak, Peon, Mali & Beat-level patrolling, fire line maintenance, nursery labour, camp duty, physical boundary marking & Frontline alert responders; use smartphones for GPS-tagged photos/voice notes; lowest formal training but highest field exposure \\
\bottomrule
\end{tabularx}
\end{table}

\textbf{Key Digital Tools:} e-Green Watch (CAMPA), AI-Based Forest Alert System, GIS/Remote Sensing platforms, Forest Offence Tracking MIS, Nursery Management System

\subsection*{Q50. Which of the following Forest department digital tools are you aware of? (Select all ...}
\qcard{Facilitating Conditions (FC)}{Department-Specific Tool Awareness}{Awareness of Forest department digital platforms}{Which of the following Forest department digital tools are you aware of? (Select all that apply)}{Multi-select Checklist (e-Green Watch / AI Forest Alert System / GIS/Remote Sensing tools / Forest Offence MIS / Nursery Management System / None)}{Forest department has a unique tech stack combining GIS, AI-based monitoring, and field verification apps; awareness must be measured at this specialised tool level to identify FC gaps.}{forbg}{forcol}

\subsection*{Q51. If your division uses the AI-based forest alert system, to what extent has it improve...}
\qcard{Performance Expectancy (PE)}{AI \& Satellite Monitoring}{Perceived usefulness of AI-based alert system for forest protection}{If your division uses the AI-based forest alert system, to what extent has it improved your ability to detect illegal activities (encroachment, felling) compared to manual patrolling?}{5-point Improvement Scale (1=No improvement to 5=Greatly improved) + Not applicable option}{MP is the first state to pilot AI-based forest alerts; measuring PE for this cutting-edge tool provides evidence on whether advanced technology translates to field-level impact (Davis, 1989).}{forbg}{forcol}

\subsection*{Q52. How difficult do you find it to use GIS-based tools (e.g., mapping forest boundaries,...}
\qcard{Effort Expectancy (EE)}{GIS Proficiency}{Perceived difficulty of using GIS and remote sensing tools}{How difficult do you find it to use GIS-based tools (e.g., mapping forest boundaries, interpreting satellite imagery, geo-tagging field reports)?}{5-point Likert (1=Very Easy to 5=Very Difficult)}{GIS requires specialised skills not covered by generic IT training; high EE for GIS indicates need for domain-specific technical capacity building unique to the Forest department (Venkatesh et al., 2003).}{forbg}{forcol}

\subsection*{Q53. Do you have access to GPS-enabled devices (smartphones, handheld GPS units) required ...}
\qcard{Facilitating Conditions (FC)}{Field Equipment}{Availability of GPS-enabled devices and field equipment}{Do you have access to GPS-enabled devices (smartphones, handheld GPS units) required for field verification, geo-tagging, and alert response?}{Categorical (Yes, department-issued / Yes, I use my personal device / No, not available)}{Forest field operations require GPS-enabled hardware for alert verification and geo-tagged reporting; unlike office-based departments, FC here includes rugged field equipment, not just desktop computers (Venkatesh et al., 2003).}{forbg}{forcol}

\subsection*{Q54. How often does poor network coverage in forest areas prevent you from using digital t...}
\qcard{Effort Expectancy (EE)}{Remote Area Operations}{Challenges of using digital tools in forest/remote terrain}{How often does poor network coverage in forest areas prevent you from using digital tools (uploading reports, responding to alerts, accessing portals) during field duty?}{5-point Frequency (1=Never to 5=Always)}{Forest staff operate in terrain with minimal cellular coverage; this department-specific EE barrier cannot be captured by office-centric connectivity questions and requires offline-capable tool recommendations (Heeks, 2006).}{forbg}{forcol}


\section{\textcolor{hlcol}{Health Department}}

\textbf{Core Functions:}

\begin{table}[H]\centering\small
\begin{tabularx}{\textwidth}{L{3cm} X}
\toprule
\textbf{Core Function} & \textbf{Description} \\
\midrule
Maternal \& Child Health & Registration, tracking, and follow-up of pregnant women, newborns, and children for ANC, immunisation, and nutrition \\[3pt]
Disease Surveillance & Real-time monitoring, reporting, and response to outbreak-prone diseases (malaria, dengue, TB, COVID) via IHIP \\[3pt]
Immunisation \& Cold Chain & Vaccine procurement, storage temperature monitoring (eVIN), and session-wise immunisation tracking \\[3pt]
TB Elimination & Case notification, treatment monitoring, DBT support, and outcome tracking under the Nikshay portal \\[3pt]
Health Assurance & Ayushman Bharat PM-JAY card issuance, hospital empanelment, and cashless treatment claim processing \\
\bottomrule
\end{tabularx}
\end{table}


\textbf{Cadre Structure \& Digital Responsibilities:}

\begin{table}[H]\centering\small
\begin{tabularx}{\textwidth}{L{1.2cm} L{3.5cm} L{4.5cm} X}
\toprule
\textbf{Class} & \textbf{Designations} & \textbf{Core Work} & \textbf{Digital Role} \\
\midrule
\textbf{I} & CMHO (Chief Medical \& Health Officer), Civil Surgeon, Joint/Deputy Director, Senior Specialist & District health administration, programme monitoring, hospital management, policy compliance & HMIS dashboard review, Nikshay district-level monitoring, Ayushman Bharat claim oversight; supervisory digital role \\[4pt]
\textbf{II} & Block Medical Officer (BMO), Medical Officer (MO), District Programme Manager (NHM) & PHC/CHC clinical duties, block-level programme implementation, data review and reporting & HMIS data validation, Nikshay case notifications, IHIP disease reporting, ANMOL supervision; regular portal use \\[4pt]
\textbf{III} & Staff Nurse, ANM (Auxiliary Nurse Midwife), Lab Technician, Pharmacist, Health Worker, Data Entry Operator & Patient care, immunisation delivery, ANC registration, cold chain maintenance, sample collection, data entry & ANMOL MP primary users (ANMs), eVIN temperature logging, HMIS facility-level data entry; heaviest routine digital workload \\[4pt]
\textbf{IV} & Ward Boy/Attendant, Peon, Sweeper, Security Guard, Ambulance Helper & Patient handling, hospital sanitation, equipment movement, security, ambulance support & Minimal formal digital role; may interact with biometric attendance systems or assist with patient registration \\
\bottomrule
\end{tabularx}
\end{table}

\textbf{Key Digital Tools:} ANMOL MP (ANM Online), HMIS, Nikshay (TB Portal), eVIN (Vaccine Cold Chain), IHIP (Disease Surveillance), Ayushman Bharat/ABHA, MPCDSR

\subsection*{Q55. Which of the following Health department digital tools are you aware of? (Select all ...}
\qcard{Facilitating Conditions (FC)}{Department-Specific Tool Awareness}{Awareness of Health department digital platforms}{Which of the following Health department digital tools are you aware of? (Select all that apply)}{Multi-select Checklist (ANMOL MP / HMIS / Nikshay / eVIN / IHIP / Ayushman Bharat-ABHA / MPCDSR / None)}{Health operates the most disease- and programme-specific portals; each serves a distinct public health function, and awareness gaps in any single tool can have patient safety implications.}{hlbg}{hlcol}

\subsection*{Q56. To what extent has the use of ANMOL MP or ABHA Health IDs improved your ability to tr...}
\qcard{Performance Expectancy (PE)}{Patient Tracking \& Continuity}{Perceived usefulness of ANMOL/ABHA for patient record management}{To what extent has the use of ANMOL MP or ABHA Health IDs improved your ability to track and follow up with patients (pregnant women, infants, TB patients) across visits?}{5-point Improvement Scale (1=No improvement to 5=Greatly improved)}{Continuity of care is the core PE outcome for Health; digital patient tracking should reduce loss-to-follow-up, and measuring this tests whether PE translates to actual health service improvement (Davis, 1989).}{hlbg}{hlcol}

\subsection*{Q57. How often are you required to enter the same patient or programme data in both paper ...}
\qcard{Effort Expectancy (EE)}{Dual Documentation Burden}{Extent of duplicate data entry across paper registers and digital systems}{How often are you required to enter the same patient or programme data in both paper registers AND digital systems (e.g., HMIS, ANMOL)?}{5-point Frequency (1=Never to 5=Always)}{Health workers frequently maintain parallel paper and digital records due to incomplete digital trust; this dual burden is a Health-specific EE barrier that increases workload without proportional benefit (Heeks, 2003).}{hlbg}{hlcol}

\subsection*{Q58. How reliably does the eVIN system reflect the actual vaccine stock and cold chain tem...}
\qcard{Facilitating Conditions (FC)}{Cold Chain \& Supply Monitoring}{Usage and reliability of eVIN for vaccine stock management}{How reliably does the eVIN system reflect the actual vaccine stock and cold chain temperature status at your facility?}{5-point Likert (1=Very Unreliable to 5=Very Reliable)}{eVIN accuracy directly affects vaccine wastage and availability; unreliable data in eVIN indicates FC failure in a system where inaccuracy has direct patient health consequences (United Nations, 2024).}{hlbg}{hlcol}

\subsection*{Q59. To what extent do mandatory real-time disease reporting requirements (through IHIP) a...}
\qcard{Social Influence (SI)}{Real-Time Surveillance Pressure}{Impact of IHIP disease surveillance reporting requirements on workflow}{To what extent do mandatory real-time disease reporting requirements (through IHIP) add to your daily workload?}{5-point Likert (1=No additional burden to 5=Significant additional burden)}{IHIP mandates real-time outbreak reporting, creating compliance-driven SI; testing whether this is perceived as productive surveillance or bureaucratic burden reveals whether SI enhances or undermines adoption (Venkatesh et al., 2003).}{hlbg}{hlcol}


\newpage

\newpage
\section{Cadre-Specific Questions (All Departments)}

These 8 questions (Q60--Q67) are administered to \textbf{every respondent} regardless of department. They capture cadre-specific digital experiences that vary by seniority and role type.

\subsection*{Q60. What is your primary interaction with digital tools in your role?}
\qcard{Facilitating Conditions (FC)}{Role-Based Digital Access}{Nature of digital tool interaction by cadre}{In your current role, what is your primary interaction with digital tools?}{Single-select (I enter/create data in portals / I review dashboards and approve / I use tools for field verification / I do not directly use digital tools / Other)}{Different cadres interact with technology at fundamentally different levels --- data entry (Class III), supervisory review (Class I-II), field use (Class III-IV), or no interaction (Class IV). Capturing this prevents treating all ``users'' as equivalent (Venkatesh et al., 2003).}{sb}{navy}

\subsection*{Q61. Has anyone in your office been assigned ``digital duty'' --- doing portal w...}
\qcard{Facilitating Conditions (FC)}{Informal Digital Labour Division}{Whether digital work is concentrated on specific individuals}{In your office, is there a specific person (e.g., a Computer Operator or junior staff member) who does most of the digital/portal work on behalf of others?}{Categorical (Yes, one designated person / Yes, a few people share it / No, everyone does their own / Not applicable)}{In many offices, a single Class III employee handles all digital work for the unit, creating a single point of failure. This ``proxy use'' pattern means official adoption numbers overstate actual individual capability (Heeks, 2003).}{sb}{navy}

\subsection*{Q62. Do you feel that employees at your level receive sufficient training co...}
\qcard{Training \& Capacity Building}{Cadre-Appropriate Training}{Perceived adequacy of training for respondent's specific cadre}{Do you feel that the digital skills training provided is appropriate for the work expected at your level (Class I/II/III/IV)?}{5-point Likert (1=Not at all appropriate to 5=Highly appropriate)}{A Class IV Forest Guard needs GPS/photo skills; a Class I Collector needs dashboard literacy. Generic ``computer training'' fails both. This item tests whether training is calibrated to actual role needs (Government of India, 2020).}{sb}{navy}

\subsection*{Q63. If a digital system breaks down or produces an error, what do you typi...}
\qcard{Facilitating Conditions (FC)}{Error Recovery by Cadre}{Coping strategy when digital systems fail}{When a digital tool or portal produces an error or stops working during your task, what do you typically do?}{Single-select (Wait for IT support / Ask a colleague for help / Switch to paper/manual process / Try to fix it myself / Escalate to supervisor / Abandon the task)}{Error recovery behaviour varies dramatically by cadre: Class I officers escalate, Class III staff switch to paper, Class IV staff abandon. This reveals the real-world FC gap between infrastructure availability and usability (Venkatesh et al., 2003).}{sb}{navy}

\subsection*{Q64. In your experience, do senior officers (Class I/II) use digital tools ...}
\qcard{Social Influence (SI)}{Leadership Digital Modelling}{Whether senior officers visibly use digital tools themselves}{In your observation, do senior officers (Class I/II) in your department actively use digital tools themselves, or do they rely on subordinates to operate portals on their behalf?}{Categorical (They use tools themselves / They mostly rely on subordinates / Mixed --- depends on the officer / I don't know)}{SI theory predicts that visible technology use by superiors accelerates adoption among subordinates. If Class I-II officers delegate all digital work, it signals that technology is ``clerical work,'' undermining adoption motivation for all cadres (Venkatesh et al., 2003).}{sb}{navy}

\subsection*{Q65. Has the introduction of digital tools changed the kind of work expecte...}
\qcard{Performance Expectancy (PE)}{Role Evolution}{Whether digital tools have changed job role expectations}{Has the introduction of digital tools changed the nature of work expected from employees at your level?}{5-point Likert (1=No change at all to 5=Completely changed my role) + Follow-up: How? (Open-ended)}{Digitisation often adds portal management to existing duties without reducing old tasks. For Class III staff (Patwaris, ANMs, GRS), this can mean doing the same field work PLUS hours of data entry --- a net negative PE outcome (Heeks, 2006).}{sb}{navy}

\subsection*{Q66. Are there tasks that your cadre is expected to do digitally but you fe...}
\qcard{Effort Expectancy (EE)}{Cadre-Skill Mismatch}{Tasks expected digitally but perceived as beyond current skill level}{Are there any digital tasks you are expected to perform as part of your duties that you feel are beyond your current skill level?}{Categorical (Yes / No) + If yes: Which tasks? (Open-ended)}{This directly measures the EE gap between institutional expectations and individual capability at each cadre level. High ``Yes'' rates in specific cadres pinpoint where targeted upskilling is needed (Venkatesh et al., 2003).}{sb}{navy}

\subsection*{Q67. In your opinion, should digital training be different for different lev...}
\qcard{Training \& Capacity Building}{Training Differentiation}{Opinion on whether training should be cadre-differentiated}{Do you think digital skills training should be designed differently for different levels of employees (Class I/II vs. III vs. IV)?}{Categorical (Yes, definitely / Somewhat / No, same training for all) + Why? (Open-ended)}{Captures ground-level demand for cadre-differentiated training design. If respondents across cadres overwhelmingly say ``yes,'' it provides direct evidence for recommending tiered training programmes in the policy brief.}{sb}{navy}


\section{Pre-Fieldwork: Questions for Your Manager}

Before beginning data collection, clarify the following with your supervisor/manager. These questions are organised by category and are designed to prevent scope misalignment, logistical delays, and institutional friction during your 3-month timeline.

\subsection{Scope \& Expectations}
\begin{table}[H]\centering\small
\begin{tabularx}{\textwidth}{C{0.8cm} X L{5cm}}
\toprule
\textbf{\#} & \textbf{Question} & \textbf{Why This Matters} \\
\midrule
1 & What is the exact scope of ``digital tools'' for this study --- only government-mandated platforms, or also informal tools (WhatsApp, personal email)? & Prevents scope creep; defines what counts as ``digital adoption'' \\[6pt]
2 & Are there any departments or offices that are off-limits or politically sensitive for data collection? & Avoids institutional friction early \\[6pt]
3 & What is the primary deliverable you expect --- a research report, a policy brief, a presentation, or all three? & Aligns output format with institutional expectations \\[6pt]
4 & Should findings be reported at department level, district level, or both? & Determines the level of disaggregation in analysis \\
\bottomrule
\end{tabularx}
\end{table}

\subsection{Access \& Permissions}
\begin{table}[H]\centering\small
\begin{tabularx}{\textwidth}{C{0.8cm} X L{5cm}}
\toprule
\textbf{\#} & \textbf{Question} & \textbf{Why This Matters} \\
\midrule
5 & Will AIGGPA provide an official letter of introduction / authorisation for me to visit department offices? & Government employees may not cooperate without formal institutional backing \\[6pt]
6 & Who is my point of contact in each department for scheduling interviews and office visits? & Prevents wasted time trying to navigate bureaucratic hierarchies independently \\[6pt]
7 & Do I need separate ethical clearance or IRB approval, or does the AIGGPA internship authorisation suffice? & Ensures data collection is institutionally legitimate \\
\bottomrule
\end{tabularx}
\end{table}

\subsection{Logistics \& Timeline}
\begin{table}[H]\centering\small
\begin{tabularx}{\textwidth}{C{0.8cm} X L{5cm}}
\toprule
\textbf{\#} & \textbf{Question} & \textbf{Why This Matters} \\
\midrule
8 & What is the realistic timeline for completing fieldwork --- is the full 3 months available, or are there institutional deadlines (reviews, presentations) that shorten it? & Allows backward-planning from hard deadlines \\[6pt]
9 & Is there a budget for travel to District Offices, or should data collection be limited to Bhopal-based offices? & Determines whether the HO vs.\ DO comparison is feasible \\[6pt]
10 & How frequently should I provide progress updates --- weekly, fortnightly, or milestone-based? & Sets reporting cadence and prevents surprise reviews \\[6pt]
11 & Can I audio-record interviews (with consent), or is there an institutional restriction? & Determines whether verbatim transcription is possible for thematic analysis \\
\bottomrule
\end{tabularx}
\end{table}

\subsection{Data \& Confidentiality}
\begin{table}[H]\centering\small
\begin{tabularx}{\textwidth}{C{0.8cm} X L{5cm}}
\toprule
\textbf{\#} & \textbf{Question} & \textbf{Why This Matters} \\
\midrule
12 & Should individual respondents be anonymised in the final report, or can designations be mentioned? & Determines the anonymisation protocol for qualitative quotes \\[6pt]
13 & Can I access any existing departmental data (IT inventory, training records, helpdesk logs) to supplement survey findings? & Secondary data strengthens triangulation without adding respondent burden \\[6pt]
14 & Who owns the final data and report --- AIGGPA, the intern, or shared? & Clarifies publication and reuse rights \\
\bottomrule
\end{tabularx}
\end{table}

\subsection{Recommendations \& Impact}
\begin{table}[H]\centering\small
\begin{tabularx}{\textwidth}{C{0.8cm} X L{5cm}}
\toprule
\textbf{\#} & \textbf{Question} & \textbf{Why This Matters} \\
\midrule
15 & How actionable should the recommendations be --- broad policy suggestions, or specific implementable steps with cost estimates? & Calibrates the depth of Objective 4 deliverables \\[6pt]
16 & Is there an existing departmental digital strategy or IT roadmap that my recommendations should align with? & Prevents recommending things already planned or rejected \\[6pt]
17 & Will there be an opportunity to present findings to department heads or senior officials? & Shapes the format and tone of the final deliverable \\
\bottomrule
\end{tabularx}
\end{table}

\newpage
\section*{References}
\addcontentsline{toc}{section}{References}
\begin{enumerate}[leftmargin=1.5em,itemsep=3pt,label={[\arabic*]}]
\item Cordella, A., \& Bonina, C. M. (2012). A public value perspective for ICT enabled public sector reforms. \textit{Government Information Quarterly}, 29(4), 512--520.
\item Davis, F. D. (1989). Perceived usefulness, perceived ease of use, and user acceptance of information technology. \textit{MIS Quarterly}, 13(3), 319--340.
\item Government of India, Dept.\ of Personnel \& Training. (2020). \textit{Mission Karmayogi}.
\item Heeks, R. (2003). Most eGovernment-for-Development Projects Fail. University of Manchester.
\item Heeks, R. (2006). \textit{Implementing and Managing eGovernment}. Sage Publications.
\item Ndou, V. (2004). E-government for developing countries. \textit{EJISDC}, 18(1), 1--24.
\item United Nations. (2024). \textit{E-Government Survey 2024}. UN DESA.
\item Venkatesh, V., Morris, M. G., Davis, G. B., \& Davis, F. D. (2003). User acceptance of information technology. \textit{MIS Quarterly}, 27(3), 425--478.
\item World Bank. (2016). \textit{World Development Report 2016: Digital Dividends}.
\end{enumerate}

\newpage
\section*{Appendix: Understanding Operational Indicators}
\addcontentsline{toc}{section}{Appendix: Understanding Operational Indicators}

An \textbf{operational indicator} is the bridge between abstract theory and concrete measurement. It answers the question: \textit{``What specific, observable thing do I need to measure to assess this concept?''}

\subsection*{Why Indicators Matter}

\begin{table}[H]\centering\small
\begin{tabularx}{\textwidth}{L{3cm} X}
\toprule
\textbf{Without Indicators} & \textbf{With Indicators} \\
\midrule
``Measure digital adoption'' (vague --- what counts as adoption?) & ``Percentage of daily tasks completed using digital tools vs.\ paper'' (specific, countable) \\[4pt]
``Assess training quality'' (how?) & ``Respondent's self-rated confidence in using department-specific portals after training'' (measurable on a scale) \\[4pt]
``Check infrastructure'' (check what?) & ``Number of functioning, dedicated digital devices per employee at their workstation'' (observable, countable) \\
\bottomrule
\end{tabularx}
\end{table}

\subsection*{The Indicator Design Chain}

Every indicator in this questionnaire follows a four-step logic:

\begin{enumerate}[leftmargin=1.5em, itemsep=6pt]
    \item \textbf{Construct} (abstract): What theoretical concept are we testing? \\
    \textit{Example: Facilitating Conditions (FC)}
    \item \textbf{Dimension} (narrower): What specific aspect of the construct? \\
    \textit{Example: Hardware Infrastructure}
    \item \textbf{Indicator} (observable): What exact thing can we see, count, or ask about? \\
    \textit{Example: Number of dedicated digital devices available per employee}
    \item \textbf{Survey Item} (question): How do we ask the respondent about it? \\
    \textit{Example: ``What digital devices are available at your personal workstation?''}
\end{enumerate}

\subsection*{Types of Indicators Used in This Study}

\begin{table}[H]\centering\small
\begin{tabularx}{\textwidth}{L{2.5cm} L{3cm} X L{3cm}}
\toprule
\textbf{Indicator Type} & \textbf{What It Measures} & \textbf{Example from This Study} & \textbf{Data It Produces} \\
\midrule
\textbf{Awareness} & Whether the respondent knows a tool exists & ``Which digital tools are you aware of?'' & Checklist (nominal) \\[4pt]
\textbf{Usage Frequency} & How often a tool is actually used & ``How often do you use digital tools?'' & Likert / ordinal \\[4pt]
\textbf{Perception} & Subjective assessment of usefulness, difficulty, or quality & ``How difficult do you find RCMS?'' & Likert (1--5 scale) \\[4pt]
\textbf{Behaviour} & What the respondent actually does & ``What do you do when a portal produces an error?'' & Categorical \\[4pt]
\textbf{Count / Quantity} & Measurable numeric values & ``How many training sessions attended?'' & Numeric (ratio) \\[4pt]
\textbf{Ratio / Proportion} & Relative comparison & ``Percentage of tasks done digitally vs.\ paper'' & Categorical bands \\[4pt]
\textbf{Open-Ended} & Rich qualitative data & ``What would improve digital tool adoption?'' & Text (for thematic analysis) \\
\bottomrule
\end{tabularx}
\end{table}

\subsection*{Good Indicators vs.\ Bad Indicators}

\begin{table}[H]\centering\small
\begin{tabularx}{\textwidth}{C{0.8cm} L{5.5cm} L{5.5cm} L{3cm}}
\toprule
\textbf{\#} & \textbf{Bad Indicator} & \textbf{Good Indicator} & \textbf{Why} \\
\midrule
1 & ``Digital readiness of the department'' & ``Percentage of employees with a dedicated device at their desk'' & Observable, countable \\[4pt]
2 & ``Quality of internet'' & ``Frequency of internet outages per week lasting over 30 minutes'' & Specific, time-bound \\[4pt]
3 & ``Training effectiveness'' & ``Self-rated confidence in using [specific portal] before vs.\ after training'' & Measurable, comparable \\[4pt]
4 & ``Employee satisfaction with technology'' & ``Agreement that digital tools make my work faster (1--5 Likert)'' & Single-barrelled, scalable \\[4pt]
5 & ``Use of AI tools'' & ``Frequency of responding to AI-generated forest alerts per month'' & Specific action, countable \\
\bottomrule
\end{tabularx}
\end{table}

\subsection*{How Indicators Connect to Analysis}

\begin{table}[H]\centering\small
\begin{tabularx}{\textwidth}{L{3cm} L{3cm} X}
\toprule
\textbf{Indicator Type} & \textbf{Analysis Method} & \textbf{What You Can Conclude} \\
\midrule
Likert (ordinal) & Median, Mann-Whitney U, Kruskal-Wallis, Cronbach's $\alpha$ & Compare perceptions across groups (departments, cadres, age groups) \\[4pt]
Categorical (nominal) & Frequency tables, chi-square, cross-tabulation & Test associations (e.g., ``Is cadre associated with error recovery strategy?'') \\[4pt]
Numeric (ratio) & Mean, SD, t-test, ANOVA, correlation & Measure exact quantities and test differences \\[4pt]
Open-ended (text) & Thematic analysis (Braun \& Clarke, 2006) & Identify patterns, generate recommendations, contextualise quantitative findings \\
\bottomrule
\end{tabularx}
\end{table}

\end{document}
